\documentclass[12pt,a4paper]{article}

\usepackage[T1]{fontenc}
\usepackage[utf8]{inputenc}
\usepackage[romanian]{babel}
\usepackage{geometry}

\geometry{margin=2.5cm}

\title{Tutoriat 4\\Modele statistice\\Recapitulare si exercitii}
\author{}
\date{}

\begin{document}

\maketitle

\section{Recapitulare}

În acest tutoriat au fost prezentate principalele concepte utilizate pentru modelarea și gestionarea incertitudinii în Inteligența Artificială.

\subsection{Variabile aleatoare}

Variabilele aleatoare permit reprezentarea fenomenelor incerte.

Exemple:

\begin{itemize}
\item Vreme;
\item Boală;
\item Trafic;
\item Rezultatul unui test medical.
\end{itemize}

Fiecărei valori posibile îi este asociată o probabilitate.

\subsection{Repartiții de probabilitate}

O repartiție de probabilitate descrie probabilitățile tuturor valorilor unei variabile aleatoare.

Proprietăți:

\begin{itemize}
\item probabilitățile sunt cuprinse între 0 și 1;
\item suma probabilităților este egală cu 1.
\end{itemize}

\subsection{Repartiții marginale și condiționate}

Repartițiile marginale se obțin prin eliminarea unor variabile dintr-o distribuție comună.

\[
P(X)
=
\sum_y P(X,y)
\]

Probabilitățile condiționate exprimă probabilități calculate în prezența unor informații suplimentare:

\[
P(A|B)
\]

\subsection{Inferența probabilistă}

Inferența probabilistă reprezintă procesul de calcul al probabilităților necunoscute pe baza informațiilor disponibile.

Metode studiate:

\begin{itemize}
\item inferența prin enumerare;
\item formula de înmulțire;
\item regula de înlănțuire;
\item regula lui Bayes.
\end{itemize}

\subsection{Formula de înmulțire}

\[
P(A \cap B)
=
P(A|B)\cdot P(B)
\]

Aceasta permite calcularea probabilităților comune pornind de la probabilități condiționate.

\subsection{Regula de înlănțuire}

Pentru mai multe variabile:

\[
P(X_1,\dots,X_n)
=
\prod_{i=1}^{n}
P(X_i|X_1,\dots,X_{i-1})
\]

Regula permite descompunerea distribuțiilor complexe în factori mai simpli.

\subsection{Regula lui Bayes}

\[
P(A|B)
=
\frac
{P(B|A)\cdot P(A)}
{P(B)}
\]

Aceasta permite actualizarea probabilităților pe baza unor observații noi.

\subsection{Independența}

Două evenimente sunt independente dacă:

\[
P(A|B)
=
P(A)
\]

În acest caz, cunoașterea lui B nu oferă informații suplimentare despre A.

\subsection{Independența condiționată}

Două evenimente sunt independente condiționat de un al treilea eveniment dacă:

\[
P(A|B,C)
=
P(A|C)
\]

Această proprietate stă la baza Rețelelor Bayesiene.

\subsection{Rețele Bayesiene}

O Rețea Bayesiană este un graf orientat aciclic în care:

\begin{itemize}
\item nodurile reprezintă variabile aleatoare;
\item muchiile reprezintă dependențe;
\item fiecărui nod îi este asociat un tabel de probabilități condiționate.
\end{itemize}

Avantaje:

\begin{itemize}
\item reprezintă eficient incertitudinea;
\item permit inferență probabilistă;
\item reduc numărul de probabilități care trebuie memorate;
\item exploatează independența condiționată.
\end{itemize}

\subsection{Idei de reținut}

\begin{enumerate}
\item Probabilitățile sunt utilizate pentru gestionarea incertitudinii.
\item Variabilele aleatoare reprezintă fenomene incerte.
\item Repartițiile descriu probabilitățile valorilor unei variabile.
\item Probabilitățile condiționate folosesc informații suplimentare.
\item Inferența probabilistă calculează probabilități necunoscute.
\item Regula lui Bayes actualizează probabilitățile pe baza evidențelor observate.
\item Independența condiționată permite simplificarea modelelor probabilistice.
\item Rețelele Bayesiene reprezintă unul dintre cele mai importante modele probabilistice din Inteligența Artificială.
\end{enumerate}

\section{Exerciții}

\subsection{Exercițiul 1 - Variabile aleatoare}

Indicați care dintre următoarele reprezintă variabile aleatoare:

\begin{enumerate}
\item Vremea de mâine.
\item Numărul de pe un zar.
\item Un program C++.
\item Rezultatul unui test medical.
\end{enumerate}

\subsection{Exercițiul 2 - Repartiție de probabilitate}

Fie variabila:

\[
X \in \{A,B,C\}
\]

cu probabilitățile:

\[
P(A)=0.2
\]

\[
P(B)=0.3
\]

\[
P(C)=0.5
\]

Verificați dacă distribuția este validă.

\subsection{Exercițiul 3 - Repartiție marginală}

Considerăm distribuția comună:

\begin{center}
\begin{tabular}{|c|c|c|}
\hline
Plouă & Trafic & Probabilitate \\
\hline
Da & Mare & 0.4 \\
Da & Mic & 0.1 \\
Nu & Mare & 0.2 \\
Nu & Mic & 0.3 \\
\hline
\end{tabular}
\end{center}

Calculați:

\[
P(Ploua=Da)
\]

și

\[
P(Ploua=Nu)
\]

\subsection{Exercițiul 4 - Probabilitate condiționată}

Se cunosc:

\[
P(A \cap B)=0.24
\]

\[
P(B)=0.3
\]

Calculați:

\[
P(A|B)
\]

\subsection{Exercițiul 5 - Formula de înmulțire}

Se cunosc:

\[
P(Ploaie)=0.3
\]

\[
P(Trafic=Mare|Ploaie)=0.8
\]

Calculați:

\[
P(Trafic=Mare \cap Ploaie)
\]

\subsection{Exercițiul 6 - Regula lui Bayes}

Se cunosc:

\[
P(Boala)=0.02
\]

\[
P(TestPozitiv|Boala)=0.9
\]

\[
P(TestPozitiv)=0.05
\]

Calculați:

\[
P(Boala|TestPozitiv)
\]

\subsection{Exercițiul 7 - Independență}

Se cunosc:

\[
P(A)=0.5
\]

\[
P(B)=0.4
\]

\[
P(A \cap B)=0.2
\]

Determinați dacă A și B sunt independente.

\subsection{Exercițiul 8 - Independență condiționată}

Explicați diferența dintre:

\begin{itemize}
\item independență;
\item independență condiționată.
\end{itemize}

Dați un exemplu pentru fiecare situație.

\subsection{Exercițiul 9 - Rețele Bayesiene}

Considerăm rețeaua:

\begin{verbatim}
Ploaie --> Trafic --> Intarziere
\end{verbatim}

Identificați:

\begin{itemize}
\item nodurile;
\item muchiile;
\item părintele variabilei Trafic;
\item copilul variabilei Trafic.
\end{itemize}

\subsection{Exercițiul 10 - Recapitulare}

Completați spațiile libere:

\begin{enumerate}
\item O variabilă aleatoare poate lua mai multe ............ .
\item Suma probabilităților dintr-o distribuție este egală cu ............ .
\item Probabilitatea condiționată se notează ............ .
\item Formula lui Bayes permite actualizarea ............ .
\item Rețelele Bayesiene sunt grafuri orientate ............ .
\item Independența condiționată stă la baza ............ .
\end{enumerate}

\section{Răspunsuri exerciții}

\begin{enumerate}

\item Variabile aleatoare:
Vremea, numărul de pe un zar, rezultatul unui test medical.

\item Da, distribuția este validă deoarece:
0.2 + 0.3 + 0.5 = 1.

\item
P(Ploua=Da)=0.4+0.1=0.5

P(Ploua=Nu)=0.2+0.3=0.5

\item
P(A|B)=0.24/0.3=0.8

\item
$P(Trafic=Mare \cap Ploaie)=0.8 \cdot 0.3=0.24$

\item
P(Boala|TestPozitiv)=$\frac{0.9 \cdot 0.02}{0.05}$=0.36

\item
$P(A)\cdot P(B)=0.5\cdot0.4=0.2$

Cum:

$P(A\cap B)=0.2=0.5\cdot0.4$

rezultă că A și B sunt independente.

\item

Independență:
\[
P(A|B)=P(A)
\]

Independență condiționată:
\[
P(A|B,C)=P(A|C)
\]

În primul caz evenimentele sunt independente în mod absolut, iar în al doilea devin independente după observarea lui C.

\item

Noduri:
Ploaie, Trafic, Întârziere

Muchii:
Ploaie \\rightarrow Trafic,
Trafic \\rightarrow Întârziere

Părinte(Trafic):
Ploaie

Copil(Trafic):
Întârziere

\item

\begin{enumerate}
\item valori;
\item 1;
\item P(A|B);
\item probabilităților;
\item aciclice;
\item Rețelelor Bayesiene.
\end{enumerate}

\end{enumerate}

\end{document}